\documentclass[aop,preprint]{../aop-template/imsart}
\RequirePackage{amsthm,amsmath,amsfonts,amssymb}
\RequirePackage[numbers]{natbib}

\startlocaldefs
\newcommand{\PTX}{\mathrm{PTX}}
\endlocaldefs

\begin{document}

\begin{frontmatter}
\title{The proof boundary for NVIDIA execution}
\runtitle{The proof boundary for NVIDIA execution}

\begin{aug}
\author{\fnms{Anonymous}~\snm{Author}}
\end{aug}

\begin{abstract}
Roop can generate and launch PTX on NVIDIA hardware. A runtime test passed on an NVIDIA GeForce RTX 4090 with driver 595.91.07, covering kernel launch and writable-buffer copyback. This is execution evidence, not a proof that CUDA execution refines Roop Lean semantics. We state the missing source, compiler, and backend claims, then give a narrow route from an integer additive kernel to a refinement theorem. The separation prevents a hardware regression test from being overstated as compiler verification.
\end{abstract}

\begin{keyword}[class=MSC]
\kwdgroup[type=primary]{\kwd{68N20}}
\end{keyword}

\begin{keyword}
\kwd{CUDA}
\kwd{compiler verification}
\kwd{reversible computation}
\end{keyword}

\end{frontmatter}

\section{Three claims}

Let $P$ be a well-formed Roop program, $\operatorname{Sem}(P)$ its Lean state transition, and $K(P)$ the emitted kernel. The source theorem has the form
\begin{equation}
\operatorname{Sem}(\operatorname{inv}(P))(\operatorname{Sem}(P)(s))=s,
\end{equation}
with the language guards, stack state, and supported effects included in the model.

A compiler theorem requires a correspondence relation
\begin{equation}
K(P)\sqsubseteq\operatorname{Sem}(P).
\end{equation}
A backend theorem relates a PTX launch and device memory to the kernel semantics:
\begin{equation}
\operatorname{launch}_{\PTX}(K(P),m)=m'
\implies m'\sim\operatorname{Sem}(P)(m).
\end{equation}
Only the conjunction of these statements justifies a claim about the running GPU program.

\section{Current evidence}

The CUDA regression test passed on the RTX 4090 configuration recorded in the repository benchmark record. It tests PTX launch and host-visible writable-buffer copyback. It does not cover arbitrary kernels, block shapes, alias patterns, memory ordering, launch failures, or floating-point behavior.

The repository contains no CUDA throughput comparison. It would therefore be inaccurate to call the backend competitive. A performance report must name the device, driver, CUDA toolchain, tensor shape, element type, operation count, launch configuration, warm-up policy, repetitions, variance, data-residency model, and baseline.

\section{A narrow theorem}

Begin with an integer additive kernel over `BitVec 64` buffers. Give it the functional specification
\begin{equation}
p\leftarrow p+f(q),
\end{equation}
where $q$ remains unchanged. Prove that the compiler emits the kernel only when its index and alias conditions hold. Then validate or prove a PTX refinement for that restricted launch form. The inverse kernel must compute $p\leftarrow p-f(q)$ from the same preserved $q$.

The certificate calculus provides the matching source-level theorem. It proves exact modular round trips but does not identify a generated kernel with its specification. Hardware tests remain useful regression tests. They do not close the compiler or backend proof boundary.

\end{document}
